% EnergyMind Tunisia -- Goal 4: integration-ready section fragment
% This file intentionally contains no \documentclass, preamble, title page,
% table of contents, document environment, or bibliography.
% Packages normally needed in the main report: amsmath, booktabs, tabularx, array.
% Citation keys used here must exist in the main report's bibliography:
% hart1992, sweeney2002, mcmahan2017, yin2018, bonawitz2017,
% dwork2014, abadi2016, mironov2017, bagdasaryan2020.

\section{Shared Learning Across the Noosphere}
\label{sec:goal4-shared}

\subsection{Objective and scope}
Goal~4 extends the local trust layer developed for Goals~1--3 into a privacy-aware collaboration between independent energy nodes. A node may represent a smart-metered neighbourhood, factory, or microgrid. Nodes exchange compact evidence about anomaly patterns, calibration behaviour, and detection knowledge so that one node can benefit from another's experience without transferring raw customer readings.

The distinction between implementation and design is important. The project includes a single-node trust-layer prototype and an in-process simulation of three nodes. The simulation is not a deployed distributed service: network transport, node authentication, and secure aggregation remain design requirements. The multi-node experiment must be run on the complete dataset before any federated-performance improvement is reported.

\subsection{Dataset and participating nodes}
The evaluation dataset is synthetic: 1,051,200 half-hourly meter rows covering 2025, generated with seed 2025 for 60 meters connected to seven feeders. It contains three simulated neighbourhoods. Public Tunisian figures informed some baseline assumptions, but the data are not measurements from a real Tunisian grid. The dataset and anomaly-generation process are described in the main data section; the relevant node partition is summarised here.

\begin{table}[htbp]
\centering
\small
\caption{Simulated nodes used in the Noosphere design.}
\label{tab:goal4-nodes}
\begin{tabularx}{\textwidth}{@{}l X r r@{}}
\toprule
Node & Simulated profile & Meters & Feeders \\
\midrule
N1 & Coastal suburb & 30 & 3 \\
N2 & Hot interior town & 15 & 2 \\
N3 & Mixed residential and commercial area & 15 & 2 \\
\bottomrule
\end{tabularx}
\end{table}

Each node keeps its raw dataset and its local models, including the local trust-score and cause-classification models. Node differences in climate, customer mix, meters, and feeder conditions motivate collaboration, but the three nodes are generated by the same simulator and therefore do not represent the full heterogeneity of independent real-world operators.

\subsection{Threat and privacy model}
The design protects customer consumption patterns, operator-specific information, and the integrity and availability of shared artifacts. Half-hourly readings can reveal occupancy patterns, vacations, sleeping hours, and appliance use, which makes raw smart-meter time series sensitive \cite{hart1992}.

The main threats considered are: (i) an honest-but-curious aggregator that analyses submitted updates; (ii) an eavesdropper that reads messages in transit; (iii) a peer node attempting to infer another node's behaviour from shared artifacts; and (iv) a corrupted node that submits misleading updates to introduce blind spots or trigger false alarms. Collusion by a majority of nodes or the creation of multiple fake node identities is outside the protection offered by the proposed aggregation rules and requires controlled enrolment and operational governance.

\subsection{What is shared and what remains local}
The governing principle is data minimisation: share the least detailed artifact that can provide cross-node value, and keep records that can identify a household or explain an individual decision inside the node.

\begin{table}[htbp]
\centering
\small
\caption{Shared artifacts, safeguards, and rationale.}
\label{tab:goal4-sharing}
\begin{tabularx}{\textwidth}{@{}p{2.5cm} p{4.0cm} p{3.5cm} X@{}}
\toprule
Information category & Shared artifact & Safeguard & Purpose \\
\midrule
Anomaly patterns & Node-level centroids summarising trusted and untrusted rows in a locally standardised feature space. & Clip feature values, perturb released centroids with noise, and combine node contributions using a coordinate-wise median. & Provides a compact signature of suspicious behaviour without sharing individual time series. These are binary trust signatures, not a complete library of subtype-specific fingerprints. \\
\addlinespace
Calibration insights & Aggregated median of the log ratio between measured energy and the physical estimate based on voltage, current, power factor, and interval duration, grouped by meter model. & Release a group only when at least three meters of that model contribute; merge contributions using medians. & Helps distinguish a systematic meter-model calibration offset from a fault affecting an individual meter. \\
\addlinespace
Relevant knowledge & Clipped and noised updates to a compact shared logistic-regression screening model. & Bound each node's update, add Gaussian noise, and use robust aggregation. & Combines learning across nodes, particularly where each node has few examples of a rare event. The shared model supplements rather than replaces local models. \\
\bottomrule
\end{tabularx}
\end{table}

The following information stays local: raw measurements and timestamps; meter, customer, and feeder identifiers; coordinates and dwelling attributes; per-meter baselines and local scaling statistics; individual-row features and verdicts; detailed explanations such as row-level SHAP values; operator notes and overrides; full decision logs; and local tree-based models. These records are retained locally for privacy, accountability, and operational control. Only node-level, protected summaries are eligible to leave a node.

\subsection{Sharing channels and protocol}
\paragraph{Local preprocessing.} Each node standardises features using its own median and median absolute deviation (MAD), which remain local, then clips the standardised values:
\begin{equation}
 z_{ij}=\operatorname{clip}\!\left(
 \frac{x_{ij}-\operatorname{median}_j}{1.4826\,\operatorname{MAD}_j},-c,c\right),
 \qquad c=5.
 \label{eq:goal4-clipping}
\end{equation}
The feature vector may include physical-consistency residuals, deviations from the meter's own history, deviations from feeder peers, feeder energy-balance changes, frozen-value indicators, and contextual features. Local standardisation statistics are not sent to the aggregator. In a deployed system, rows that the local trust layer has quarantined or escalated should be excluded from shared training; this quarantine-before-sharing control is a design requirement and must not be described as implemented unless it is verified in the code.

\paragraph{Federated screening model.} For round $t$, node $k$ trains locally from the current shared weights $w_t$ and returns an update. The proposed clipping-and-noise step is
\begin{align}
 \Delta_k &= \operatorname{LocalTrain}(w_t;\mathcal{D}_k)-w_t, \\
 \widetilde{\Delta}_k &= \Delta_k\min\!\left(1,\frac{C}{\lVert\Delta_k\rVert_2}\right)
 + \mathcal{N}\!\left(0,\frac{\sigma^2 C^2}{K}I\right), \\
 w_{t+1} &= w_t+\operatorname{Agg}(\widetilde{\Delta}_1,\ldots,\widetilde{\Delta}_K),
 \label{eq:goal4-fed-update}
\end{align}
where $C$ is the clipping bound, $\sigma$ is the noise multiplier, $K$ is the number of participating nodes, and $\operatorname{Agg}$ is the selected aggregation rule. The current design considers a mean (FedAvg), coordinate-wise median, or reputation-weighted aggregation \cite{mcmahan2017,yin2018}. The calibration digest and the anomaly signatures are exchanged less frequently than the model update.

\paragraph{Round protocol.} A sharing round follows these steps:
\begin{enumerate}
  \item The coordinator publishes the round identifier, protocol version, feature-schema identifier, clipping bound, noise setting, and submission deadline.
  \item Each node prepares eligible local examples, computes its protected model update, and periodically prepares its anomaly signature and calibration digest.
  \item The coordinator checks schema compatibility, message freshness, update norms, and obvious outliers; it then aggregates valid submissions and creates a versioned shared artifact.
  \item The coordinator distributes the artifact with its version and integrity metadata. In a deployment, messages must use authenticated, encrypted transport, and secure aggregation should prevent the coordinator from inspecting individual model updates \cite{bonawitz2017}.
  \item Each node evaluates the candidate artifact on a local held-out validation set (a \emph{canary}). The node adopts it only if the agreed acceptance criterion is met and its operator approves; otherwise, it keeps the prior version and records the refusal. A deployed system should support revocation and rollback.
\end{enumerate}

A weekly model-sharing cadence is a reasonable starting design because it matches the dataset's seven-day rolling context; an urgent advisory can communicate a newly confirmed attack signature sooner. Message versioning and a feature-schema identifier are needed to ensure that identically named feature positions have the same meaning across nodes. Exact class counts should not be released in a real deployment; counts should be bucketed or protected because they reveal a node's anomaly prevalence.

\subsection{Privacy protection and its limits}
The proposed privacy controls are complementary rather than interchangeable:
\begin{itemize}
  \item \textbf{Local data retention:} raw readings, identifiers, local explanations, and detailed logs do not leave the node.
  \item \textbf{Clipping:} limits how much one local update can influence an aggregation round.
  \item \textbf{Gaussian noise:} reduces information leaked through released updates and signatures; its scale must be tied to the sensitivity of the relevant release \cite{dwork2014,abadi2016}.
  \item \textbf{Group suppression:} the calibration digest is released only for meter-model groups containing at least three meters. This is a small-group suppression rule, not a formal differential-privacy guarantee \cite{sweeney2002}.
  \item \textbf{Secure aggregation:} proposed for deployment so that the coordinator receives only an aggregate of node updates, rather than each node's update individually \cite{bonawitz2017}.
\end{itemize}

For the mean-aggregation case, a simplified Rényi-DP accounting expression over $R$ rounds with full participation is
\begin{equation}
 \varepsilon(\delta)=\min_{\alpha>1}\left[
 \frac{R\alpha}{2\sigma^2}+\frac{\ln(1/\delta)}{\alpha-1}
 \right],
 \qquad \delta=10^{-5},
 \label{eq:goal4-dp-accounting}
\end{equation}
under the corresponding clipping, noise, and secure-aggregation assumptions \cite{mironov2017}. This expression must not be presented as a certified privacy guarantee for the complete system: it applies to the stated mean-aggregation setup and does not by itself account for the signatures, calibration digests, exact class counts, participation metadata, or every possible release. The median and reputation-weighted aggregators also require separate analysis. With only three nodes, noise sufficient for strong privacy may substantially reduce model utility.

Most importantly, the current in-process simulation does not implement secure aggregation or network security. Until those controls and the full privacy accounting are implemented and reviewed, the design should be described as privacy-aware, not as providing a proven end-to-end privacy guarantee.

\subsection{Defence against corrupted nodes and governance}
The collective-learning layer must assume that a participant can be faulty or malicious, just as the local trust layer assumes that meter data can be corrupted. A corrupted node might submit a sign-flipped or scaled update, train on falsified readings, replay a stale update, or attempt a targeted backdoor. Backdoor attacks are a known limitation of basic federated-learning defences \cite{bagdasaryan2020}.

The proposed defence combines:
\begin{enumerate}
  \item \textbf{Bounded influence:} clip each model update before aggregation.
  \item \textbf{Robust aggregation:} use a coordinate-wise median or another robust rule instead of relying solely on a simple mean \cite{yin2018}. With three nodes, a median can resist at most one arbitrary outlier in the idealised one-corrupted-node setting; it does not protect against a colluding majority.
  \item \textbf{Reputation and quarantine:} compare each update with the population of updates, down-weight or reject extreme submissions, and request human review when a node is repeatedly inconsistent. This mechanism is not a security proof and must be validated against stronger attacks.
  \item \textbf{Local canary and human approval:} test each incoming artifact on local validation data before adoption. Shared information can increase suspicion but must not independently trigger an operational grid action.
  \item \textbf{Auditability and rollback:} keep versions of accepted artifacts and record adoption, refusal, and rollback in the node's audit trail.
\end{enumerate}

These controls address both over-trust and under-trust. Over-trust is limited by local validation and operator approval; under-trust is made reviewable by recording refusal reasons and comparing candidate and current-model performance. Local escalation and autonomy rules from Goal~3 remain authoritative regardless of what the shared model recommends.

\subsection{Evaluation plan and implementation status}
The intended evaluation uses a chronological split for each node: January--August for training, September--October for calibration and validation, and November--December for testing. Rows without a valid trust label are excluded, and ground-truth columns (the dataset's \texttt{gt\_*} fields) are never model inputs. Results should compare: (i) local-only models, including cross-node transfer; (ii) clean federated learning; (iii) federated learning with clipping and noise; (iv) mean versus robust aggregation when one node is corrupted; and (v) anomaly signatures evaluated on a held-out node. Report ROC-AUC alongside precision, recall, and false-positive rate at a specified alert threshold.

No numerical claim about federated improvement should be made until this experiment has been run on the complete dataset and its outputs checked. The existing pipeline was previously exercised on a small toy table, which establishes basic execution only; it does not establish model performance on the full EnergyMind dataset.

\subsection{Limitations}
The proposed design has four principal limitations. First, all three nodes and all injected anomalies are synthetic, so the experiment may overfit the simulator and cannot establish real-grid performance. Second, the nodes share one generator and are less heterogeneous than independent utilities. Third, only three participants provide little redundancy for robust aggregation and make strong differential privacy expensive. Fourth, the tested corrupted-node scenario is limited and does not establish resistance to backdoors, collusion, adaptive attacks, or compromised aggregators. Before deployment, the system requires evaluation on more heterogeneous data, authenticated transport, secure aggregation, a reviewed privacy accountant for every shared artifact, stronger adversarial testing, and a documented node-enrolment and revocation process.
